\paragraph{Construction cost and repeated-run variation.}
Table~\ref{tab:encoding-costs} separates construction and search at $n=128$.
Build time includes generating clauses and inserting them into the SAT
solver; solve time measures the solver call. The process column gives a
median and observed minimum--maximum range across the three repetitions.
Each timing column is summarized separately, so its medians need not add
up. A range from three repeats is descriptive, not a confidence interval.
\begin{table}[pos=htbp]
\centering\small
\caption{Encoding-study times at $n=128$, in seconds. Completed cells
contain three valid solutions. No component times are imputed for timeouts.}
\label{tab:encoding-costs}\begin{tabular}{lrrr}
\toprule
Encoding / formulation & Build & Solve & Process [min, max] \\
\midrule
pairwise/minimal & 4.050 & 1.332 & 5.795 [5.474, 5.954] \\
pairwise/extended & 3.901 & 2.593 & 6.769 [6.749, 6.963] \\
binary/minimal & 0.566 & 5.385 & 6.218 [6.179, 6.407] \\
binary/extended & 0.578 & 2.514 & 3.362 [3.111, 3.646] \\
commander/minimal & 0.325 & 0.240 & 0.886 [0.858, 0.892] \\
commander/extended & 0.327 & 0.257 & 0.888 [0.867, 0.890] \\
product/minimal & -- & -- & TO (0/3) \\
product/extended & 0.247 & 2.220 & 2.736 [2.644, 2.830] \\
seqcounter/minimal & -- & -- & TO (0/3) \\
seqcounter/extended & 0.349 & 2.693 & 3.331 [3.185, 3.472] \\
\bottomrule
\end{tabular}

\end{table}
\FloatBarrier
For pairwise/minimal, the median construction time is 4.050 s, compared with 0.325 s for commander/minimal. The ratio is approximately 12.5 on this implementation. This connects the clause-count reduction to an observed construction saving. It does not determine search performance: product/extended has fewer clauses than commander/minimal but a longer solver call.

\paragraph{Effect of redundant column clauses.}
At $n=128$, product/minimal and seqcounter/minimal each finish zero of three runs, while their extended versions finish all three with median process times 2.736 s and 3.331 s. In each case, extended adds only 128 column ALO clauses and no variables. These explicit column disjunctions make the column occupancy requirement directly available to unit propagation; in minimal it follows from the combined row and column constraints. The measured improvement is consistent with a change in propagation or search, but timings alone do not isolate which internal events caused it. Commander receives the same additional clauses without a comparable improvement. The benefit is therefore conditional on the encoding, not a universal advantage of adding redundant clauses.
